\documentclass{article}
\input{../../lib.sty}
\title{\texttt{fn threshold\_delta}}
\author{\MichaelShoemateAuthor}
\begin{document}
\maketitle
\contrib

\section{Hoare Triple}
\subsection*{Precondition}
\subsubsection*{Compiler-verified}
The arguments are as in the Rust signature.
\subsubsection*{Caller-verified}
For a symmetric discrete noise distribution $Z$ with log-concave mass function,
\texttt{tail(t)} returns an upper bound on $\Pr[Z>t]$ or an error.

\subsection*{Postcondition}
\begin{theorem}
On success, the output upper bounds
\[
\max_{\substack{a\in\mathbb Z_{\ge0}^{n}\\
\sum_i a_i\le l_1,\ a_i\le b}}
\left(1-\prod_{i:a_i>0}F(T-a_i)\right),
\]
where $n=l_0$, $b=l_\infty$, $T=\texttt{threshold}$, and $F$ is the CDF of $Z$.
\end{theorem}

\begin{proof}
Let $s=\min(l_1,nb)$, replace $b$ by $\min(b,s)$, and set $h(a)=-\log F(T-a)$.
The cases $s=0$ and $b=1$ are immediate.
Otherwise $h$ is increasing and discretely convex since the CDF of a log-concave mass function is log-concave.
Thus a maximizer spends $s$ and has support $m\in[\lceil s/b\rceil,\min(n,s)]\cap\mathbb Z$.
Subtract one from every positive coordinate, and write
\[
s-m=q(b-1)+r,\qquad 0\le r<b-1.
\]
By weak majorization (as in \texttt{l01inf\_to\_lp}), the maximum of $\sum_{i:a_i>0}h(a_i)$ at fixed $m$ is
\[
H(m)=mh(1)+q(h(b)-h(1))+h(1+r)-h(1).
\]
Between successive breakpoints $r=0$, the expression is discretely convex:
the same formula extends to $r=b-1$ at the adjacent breakpoint.
At the breakpoints,
\[
H(m)=mh(1)+\frac{s-m}{b-1}(h(b)-h(1)),
\]
which is affine in $m$.
Hence it suffices to evaluate the interval endpoints and the first and last feasible breakpoints.
For each candidate, the probability is $1-\prod_i(1-p_i)$, where $p_i=\Pr[Z>T-a_i]$.
The function $u(x,y)=x+y-xy$ is increasing in both arguments on $[0,1]$.
The implementation evaluates $u$ exactly over rational representations of the supplied bounds and rounds upward.
Binary doubling therefore conservatively combines independent events without cancellation.
The maximum of the candidate bounds proves the result.
\end{proof}
\end{document}
